\chapter{The Method of Choice}

A person's values are not mainly revealed by what he praises. They are revealed by what he chooses.

Values answer two questions:

\begin{quote}
What is more important? What is most important?
\end{quote}

Once those questions have been answered clearly, choice becomes possible. Without them, choice becomes confusion. A person looks at several options, adds more and more conditions, fears the consequences, asks everyone for advice, and finally calls the whole experience ``being torn.''

Much of this pain is unnecessary. Choice is not mysterious. It is a skill. Like most skills, it improves when the process is made explicit, practiced often, and reviewed honestly.

\section*{Too Many Conditions}

Consider a familiar example: choosing a spouse.

Many people who remain single for a long time later discover that they probably had more than one real opportunity. At the time, however, they did not choose. If you ask why, they often say their standards were not high. They were not demanding the richest, the most beautiful, or the most brilliant partner. They only wanted ordinary conditions:

\begin{itemize}
\item not too unattractive;
\item not too short;
\item not too boring;
\item not too poor;
\item not too poorly educated.
\end{itemize}

Each condition sounds moderate. Suppose each filter leaves only the top third of available people. Five such filters do not leave one third. They leave:

\[
\left(\frac{1}{3}\right)^5
=
\frac{1}{243}
\approx 0.41\%.
\]

That is about four people in a thousand.

Now add reality. During the suitable years of life, how many people will one person know deeply enough to consider seriously? Perhaps not more than \(150\), including people who are not even candidates. And there is another filter: the other person is also choosing.

The problem was not bad luck. The problem was a filter set that silently eliminated almost everyone.

This mistake is not limited to marriage. It appears in careers, investing, hiring, entrepreneurship, product design, and daily life. A person says, ``My requirements are not high. Why is there no good option?'' The arithmetic answers calmly: because you added too many conditions.

Probability thinking reduces self-pity. It shows that the world is often not cruel. It is simply not shaped like our wishes.

\section*{The One Important Condition}

People who choose well often do something that looks almost unfairly simple:

\begin{quote}
They identify the most important condition, then keep their attention on it.
\end{quote}

This is the practical meaning of values. Values are not decorative beliefs. They are ranking systems. They decide which conditions matter and which conditions only feel comforting.

If shared growth is the most important condition in a marriage, then some other pleasant but secondary conditions must yield. If long-term demand is the most important condition in a career, then a fashionable title may matter less. If risk control is the first necessity of investment, then excitement is not a virtue. If attention is the most valuable personal resource, then many cheap pleasures are expensive.

Earlier chapters have returned to this point from different angles. Attention is more valuable than money. Demand decides price. Investment begins with avoiding ruin, not chasing thrills. Long-term thinking changes the meaning of opportunity. Each of these concepts is a value judgment: it tells us what should outrank what.

To refine one's values is to improve the quality of one's choices.

\section*{Choice Is the Main Event}

Among all the important matters in life, choice may be the most important.

This sounds exaggerated until we count the truly decisive choices. They are not infinite. A person chooses what to study, where to work, whom to marry, which city or industry to enter, what kind of people to spend time with, what business to build, what assets to buy, what risks to avoid, and what skills to cultivate.

There may be more, but not many more.

Most of daily life does not deserve heroic calculation. It is wasteful to spend first-class attention comparing three nearly identical household items, then approach a career or investment decision with a tired mind. Wisdom is not needed equally everywhere. It must be reserved for the few choices that can redirect years.

This is another reason to protect attention. The goal is not to become a miser of attention for its own sake. The goal is to have enough clean attention available when the choice is worth it.

\section*{Add Necessary Conditions}

The core method of choice can be stated in one sentence:

\begin{quote}
Add necessary conditions.
\end{quote}

The sentence is short, but each word matters.

To choose is to filter. A condition removes options. This is useful because a world with too many options is unworkable. But every condition has a cost. Add an unnecessary condition, and good options disappear. Omit a necessary condition, and bad options survive.

So the discipline is not ``add many conditions.'' The discipline is:

\begin{quote}
Consider only the necessary conditions, and do not miss any necessary condition.
\end{quote}

This is where Ockham's Razor becomes practical. ``Do not multiply entities beyond necessity'' is often discussed as a principle of explanation. It also works as a principle of decision. In choice, unnecessary conditions are not harmless. They distort the result.

The structure is simple:

\begin{center}
\begin{adjustbox}{max width=\linewidth}
\begin{tikzpicture}[node distance=0.95cm, every node/.style={font=\small}]
\node[draw, rounded corners, align=center, minimum width=2.8cm, minimum height=0.75cm] (options) {Options};
\node[draw, rounded corners, align=center, minimum width=3.0cm, minimum height=0.75cm, right=of options] (conditions) {Necessary\\conditions};
\node[draw, rounded corners, align=center, minimum width=2.8cm, minimum height=0.75cm, right=of conditions] (choice) {Choice};
\node[draw, rounded corners, align=center, minimum width=3.2cm, minimum height=0.75cm, below=of conditions] (values) {Values decide\\necessity};
\draw[->] (options) -- (conditions);
\draw[->] (conditions) -- (choice);
\draw[->] (values) -- (conditions);
\end{tikzpicture}
\end{adjustbox}
\end{center}

A poor chooser adds conditions to calm anxiety. A good chooser adds conditions to reveal reality.

\section*{Amazon's First Product}

Amazon began as an online bookstore. In hindsight, this feels obvious. Books were the first product, and the choice worked. But the usefulness of the case lies in the filtering process.

The team did not merely ask, ``What can we sell online?'' They searched for a product category with several necessary conditions:

\begin{itemize}
\item the market had to be large;
\item the category had to have long-term growth;
\item customers had to buy repeatedly;
\item after-sales cost had to be low, preferably close to zero.
\end{itemize}

Each condition removed many categories. If each condition eliminated \(90\%\) of possible products, four conditions would leave roughly one in ten thousand.

That is severe filtering. It is also exactly what a serious choice may require.

The lesson is not that every business should sell books. The lesson is that high-quality choices often look simple only after the fact. The simplicity came from disciplined filtering. The difficulty was knowing which filters were necessary.

This is why successful teams are not necessarily teams that ``thought of everything.'' No one can think of everything before facing an uncertain future. Better teams do something more realistic and more valuable: they think harder about the necessary conditions.

\section*{Not Choosing Is Also Choosing}

Many people try to escape difficult choices. They wait. They ask others to decide. They search for a perfect rule. They discuss abstract questions like whether human beings really have choice at all.

But not choosing is itself one option among the options.

The cost is that the person loses practice. He gives the decision away, then also gives away the chance to improve his own decision-making ability. The next choice arrives, and he is still weak. This becomes a loop: fear produces avoidance, avoidance prevents practice, lack of practice produces more fear.

The only way out is to treat choice as trainable.

A person who thinks seriously may look tired to others. Someone may ask, ``Isn't it exhausting to think so much?'' The better question is: what is the cost of not thinking now? A poorly made choice can create years of avoidable trouble. Thirty minutes of serious thought may prevent months or years of repair.

Thinking is not the enemy of life. Bad thinking is.

\section*{The Two Common Errors}

When reviewing bad decisions, the same two errors appear again and again:

\begin{enumerate}
\item a necessary condition was ignored;
\item a necessary condition was considered, but not strictly enough.
\end{enumerate}

This applies especially to investment. Most failed investments do not require a mysterious explanation. Somewhere in the decision process, an essential filter was missing or loose. The market was not large enough. Growth was weak. The team was not trustworthy. The risk was misunderstood. The investor wanted excitement more than safety. The necessary condition existed, but the investor did not treat it as necessary.

This is painful because it removes many comforting excuses. Luck matters. Timing matters. Uncertainty is real. Even a correct choice can fail. But over a long enough period, a person's present situation is heavily shaped by earlier choices.

The longer the time horizon, the harder it becomes to blame everything on fate.

This realization can be uncomfortable, but it is also liberating. If poor choices created many present difficulties, better choices can create a different future.

\section*{Choice Before Execution}

A useful line about entrepreneurship says:

\begin{quote}
Successful entrepreneurship is a group of people good at essay questions getting the multiple-choice question right.
\end{quote}

Execution matters. Persistence matters. Skill matters. But before all of that, the direction must be chosen. Doing the wrong thing with great discipline is still expensive.

An older formula says:

\begin{quote}
Success is doing the right thing in the right way.
\end{quote}

The right way cannot rescue the wrong thing indefinitely. This is why investors often ask founders about the basic filters of a business direction. Does it satisfy high frequency, rigid demand, and a large market? If not, why does that not matter here? The answer reveals more than the content of the answer. It reveals whether the founder has treated the choice itself seriously.

A short conversation about necessary conditions can expose a great deal. Some founders have slogans. Some have excuses. Some have real thinking. The difference matters.

\section*{A Practice Method}

Decision-making can be practiced with a notebook.

The exercise is deliberately plain:

\begin{enumerate}
\item When facing a choice, list the possible options.
\item List the conditions by which the options might be filtered.
\item Give each condition an importance score from \(1\) to \(5\).
\item Mark whether each condition is necessary: \(1\) for necessary, \(0\) for not necessary.
\item Remove unnecessary conditions, then apply the necessary ones strictly.
\item Wait until the next day and repeat the process.
\item For important choices, repeat several times before acting.
\end{enumerate}

The waiting period matters. A first pass often contains mood, fear, vanity, and borrowed opinions. The second pass is usually cleaner. For a major choice, thirty minutes today and thirty minutes tomorrow is a small price.

The notebook should not record only the final decision. It should record the process: the options, the filters, the ranking, the doubts, the rejected conditions, and the reason for acting. Without process records, later review becomes storytelling. With records, review becomes learning.

\begin{center}
\begin{adjustbox}{max width=\linewidth}
\begin{tabular}{llll}
\hline
Condition & Importance & Necessary? & Notes \\
\hline
Large demand & 5 & 1 & Must be present \\
Low ruin risk & 5 & 1 & No compromise \\
Personal interest & 3 & 0 & Helpful, not decisive \\
Fashionable label & 1 & 0 & Remove \\
\hline
\end{tabular}
\end{adjustbox}
\end{center}

The exact table is not important. The habit is important.

\section*{Values and Choices Form a Loop}

Values shape the conditions we add. The results of those choices then reshape our values.

If a person chooses by vanity, he eventually pays tuition to reality. If he records the process and reviews honestly, the pain can become education. He may discover that vanity was not a necessary condition. He may learn that safety, growth, demand, trust, or long-term compounding should have ranked higher.

This is the positive loop:

\begin{center}
\begin{adjustbox}{max width=\linewidth}
\begin{tikzpicture}[node distance=1.0cm, every node/.style={font=\small}]
\node[draw, rounded corners, align=center, minimum width=2.7cm, minimum height=0.75cm] (values) {Values};
\node[draw, rounded corners, align=center, minimum width=2.7cm, minimum height=0.75cm, right=of values] (filters) {Necessary\\conditions};
\node[draw, rounded corners, align=center, minimum width=2.7cm, minimum height=0.75cm, right=of filters] (choice) {Choice};
\node[draw, rounded corners, align=center, minimum width=2.7cm, minimum height=0.75cm, below=of filters] (review) {Review};
\draw[->] (values) -- (filters);
\draw[->] (filters) -- (choice);
\draw[->] (choice) |- (review);
\draw[->] (review) -| (values);
\end{tikzpicture}
\end{adjustbox}
\end{center}

Without records, the loop breaks. Memory is too convenient. It forgets embarrassment, edits motives, and invents reasons after the outcome is known. Many people remain unchanged not because they lack intelligence, but because they never preserve the evidence needed for correction.

A decision journal is one of the simplest tools against endless hesitation. It turns vague struggle into inspectable material.

\section*{The Danger of Being Too Clever}

Some people fail to grow because they are clever in the wrong direction. Their intelligence is spent searching for an easier method. They want the shortcut before doing the ordinary practice. They collect better frameworks, better slogans, and better advice, but do not apply the simple method already available.

Choice does not improve through admiration of decision theory. It improves through decisions made, recorded, reviewed, and corrected.

This is especially important for investors. Investment is an industry built on choice. What to buy, what not to buy, when to wait, how much to allocate, which risk to avoid, which opportunity is not worth attention: these are all choices. A person who refuses serious choice practice is not ready for serious investing.

He may still buy assets. He may still get lucky. But luck without choice ability is not freedom. It is exposure.

\section*{Questions for Major Choices}

For the important decisions of life, the following questions deserve written answers:

\begin{enumerate}
\item What am I choosing among?
\item What outcome matters most here?
\item Which conditions are truly necessary?
\item Which conditions merely reduce anxiety or satisfy vanity?
\item Which necessary condition am I tempted to relax?
\item What consequence am I unwilling to bear?
\item If this choice fails, what will probably be the reason?
\item What would make this choice look wise ten years from now?
\end{enumerate}

The question about consequence is essential. Many people do not lack options. They lack willingness to bear the result of choosing. So they remain undecided, hoping reality will choose for them. Reality usually does, but not always kindly.

Maturity includes accepting that no choice gives everything. Growth is the process of realizing that one cannot keep all possible lives, then learning to choose the one that deserves sacrifice.

\section*{The End of Complaint}

As a person's attitude toward major choices becomes more serious, complaint weakens.

In youth, it is easy to believe one's circumstances are mostly unfair. Later, with enough review, another pattern appears: many embarrassments were chosen earlier, even if unconsciously. The choice of friends, industry, habits, spouse, business model, reading, spending, risk, attention, and patience all accumulated.

This does not mean every suffering is deserved. It means a useful person should search first for the part that can be improved by better choice.

The road to wealth freedom is not paved by perfect predictions. It is built by improving the quality of repeated choices under uncertainty. Choose where to place attention. Choose what to learn. Choose what demand to serve. Choose what risk to avoid. Choose whom to build with. Choose which necessary conditions cannot be compromised.

Then record the process.

Tomorrow's values will be shaped by today's choices. Today's choices can be improved by clearer values. Keep the loop running in the right direction.
